\section*{Capítulo \ref{ch:b3:quantum-formalism} --- El formalismo de la mecánica cuántica}

\begin{solution}{exo:b3:quantum-formalism:1}
(a) $(4 + 1)/5 = 1$ y $(1 + 1)/2 = 1$. (b) $\braket\varphi\psi =
(2 - \iu)/\sqrt{10}$; $\braket\psi\varphi = (2 + \iu)/\sqrt{10}$:
conjugados. (c) $|(2 - \iu)|^2/10 = 1/2$. (d) Resuélvase
$\braket\psi\chi = 0$: $\ket\chi = (\ket1 - 2\iu\ket2)/\sqrt5$.
\end{solution}

\begin{solution}{exo:b3:quantum-formalism:2}
Primera: hermítica; $\pm1$ con $(\ket1 \pm \ket2)/\sqrt2$. Segunda:
hermítica; $\pm1$ con $(\ket1 \pm \iu\ket2)/\sqrt2$. Tercera: no
hermítica (la transpuesta conjugada difiere). Cuarta: hermítica;
$\lambda^2 - 5\lambda + 4 = 0$ da $1$ y $4$, con vectores propios
$\big({-(1 - \iu)}\ket1 + \ket2\big)/\sqrt3$ y $\big(\ket1 +
(1 + \iu)\ket2\big)/\sqrt3$.
\end{solution}

\begin{solution}{exo:b3:quantum-formalism:3}
(a) $\braket\beta\alpha = \cos\beta\cos\alpha + \sin\beta\sin\alpha
= \cos(\beta - \alpha)$: probabilidad $\cos^2(\beta - \alpha)$. (b)
La intensidad es $\propto$ al número de fotones: Malus. (c) Cada fotón es
un ensayo independiente: varianza $N\mathcal P(1 - \mathcal P)$ --- el
propio ruido certifica los fotones. (d) $|\cos\beta + \iu\sin\beta|^2/2
= 1/2$ para todo $\beta$: la luz circular no privilegia ningún eje
transversal.
\end{solution}

\begin{solution}{exo:b3:quantum-formalism:4}
(a) $(\hat x\hat p - \hat p\hat x)f = -\iu\hbar(xf' - (xf)') =
\iu\hbar f$. (b) $[\hat x^2, \hat p] = 2\iu\hbar\hat x$; $[\hat x,
\hat p^2] = 2\iu\hbar\hat p$. (c) Súmese y réstese $\hat B\hat A
\hat C$. (d) Inducción con (c): $\iu\hbar\,n\hat p^{n-1}$ ---
$[\hat x, \cdot]$ actúa como $\iu\hbar\,\partial/\partial p$, la
sombra cuántica del \omterm{def:b3:hamiltonian-mechanics:poisson}{corchete de Poisson} con $x$.
\end{solution}

\begin{solution}{exo:b3:quantum-formalism:5}
(a) $|\braket{45^\circ}{V}|^2 = \tfrac12$, colapso y luego
$|\braket{H}{45^\circ}|^2 = \tfrac12$: en total, $\tfrac14$. (b) Tomándolos
en el orden $60^\circ$, $30^\circ$ (tres pasos de $30^\circ$):
$(\cos^230^\circ)^3 = 27/64 \approx 0.42$. (c) Cada paso de
$\pi/2N$ se atraviesa con $\cos^2(\pi/2N)$: supervivencia $[\cos^2(\pi/2N)]^N
= 0.25$, $0.60$ y $0.88$ para $N = 2, 5, 20$. (d) Cuando $N \to \infty$
la supervivencia tiende a $1$: muchas medidas suaves conducen el estado a
lo largo de $90^\circ$ sin pérdida --- la medida usada como volante.
\end{solution}

\begin{solution}{exo:b3:quantum-formalism:6}
(a) $[\hat A, \hat B] = \begin{pmatrix}0 & -2\\ 2 & 0\end{pmatrix}
\neq 0$: incompatibles. (b) $\hat B$ sobre $\ket1$: $+1$ con
certeza, sin necesidad de colapso; $\hat A$ da después $\pm1$ con
probabilidad $\tfrac12$ cada uno. (c) Tras $\hat A \to +1$, el estado es
$(\ket1 + \ket2)/\sqrt2$; $\hat B$ lo colapsa a $\ket1$ o a
$\ket2$ (con $\tfrac12$ cada uno); en cualquier caso, la última medida de
$\hat A$ da $\pm1$ con probabilidad $\tfrac12$: contradicción con
probabilidad $\tfrac12$. (d) Los \omterm{def:b3:quantum-formalism:observable}{observables} que conmutan comparten estados
propios: la medida intermedia no perturbaría, y la repetición coincidiría
con certeza.
\end{solution}

\begin{solution}{exo:b3:quantum-formalism:7}
(a) $\Delta x = \sigma$ y $\Delta p = \hbar/2\sigma$ (la transformada de
Fourier de una gaussiana de anchura $\sigma$ tiene anchura $1/2\sigma$ en
$k$): el producto vale exactamente $\hbar/2$. (b) Igualdad en
Cauchy--Schwarz: $\hat\beta\ket\psi \propto \hat\alpha\ket\psi$ con razón
imaginaria pura --- para $x, p$ esa ecuación diferencial solo tiene
soluciones gaussianas. (c) $\Delta p \ge \qty{5.3e-25}{kg.m/s}$: $E \sim
\Delta p^2/2m_{\text{e}} \approx \qty{1}{eV}$ --- los átomos son máquinas
de electronvoltios porque son cajas de ángstroms. (d)
$\Delta v \ge \qty{5e-27}{m/s}$: nada, jamás.
\end{solution}

\begin{solution}{exo:b3:quantum-formalism:8}
(a) $[\hat H, \hat x] = -\iu\hbar\hat p/m$ y $[\hat H, \hat p] =
\iu\hbar V'(\hat x)$ dan $\dd\langle x\rangle/\dd t = \langle
p\rangle/m$ y $\dd\langle p\rangle/\dd t = -\langle V'\rangle$. (b)
Un cambio de variable en el producto interno muestra la hermiticidad, y
$\hat\Pi^2 = \mathbb 1$ obliga a \omterm{def:b3:quantum-formalism:observable}{valores propios} $\pm1$. (c) $V(-x) =
V(x)$ hace a $\hat H$ ciego a la paridad; un estado propio no degenerado
debe entonces ser también estado propio de $\hat\Pi$: par o impar. (d) El
doblete casi degenerado del pozo doble simétrico: un estado par y uno impar
--- la pareja de inversión del amoniaco, desdoblada por efecto túnel, que
radia a \qty{24}{GHz}.
\end{solution}

\begin{solution}{exo:b3:quantum-formalism:9}
(a) Ambos estados comparten $E$: anunciar la energía deja dos
posibilidades. (b) $\hat S$ intercambia las etiquetas: es hermítico, su
cuadrado es la identidad y conmuta con el $\hat H$ simétrico; dentro del
nivel, sus estados propios son $(\ket{1,2} \pm \ket{2,1})/\sqrt2$ con
\omterm{def:b3:quantum-formalism:observable}{valores propios} $\pm1$. (c) $(E, +)$ y $(E, -)$: etiquetas únicas. (d) Una
degeneración es una dirección incompleta; la simetría que la causa
proporciona también la cifra que falta.
\end{solution}

\begin{solution}{exo:b3:quantum-formalism:10}
(a) $\Delta E\,\Delta A \ge \tfrac12|\langle[\hat H, \hat
A]\rangle| = \tfrac\hbar2|\dd\langle A\rangle/\dd t|$; divídase. (b)
El tiempo es un parámetro de la teoría, no un operador: la relación
acota lo \emph{deprisa} que puede evolucionar cualquier magnitud medible,
dada la dispersión de energía. (c) $\Delta E \sim \hbar/2\tau = \qty{3.3e-8}{eV}$:
una anchura natural de línea de unos pocos megahercios. (d) Un cambio
visible por segundo exige solo $\Delta E \gtrsim \qty{5e-35}{J}$ --- para
energías macroscópicas, ninguna restricción: los relojes pueden marchar.
\end{solution}

\begin{solution}{exo:b3:quantum-formalism:11}
(a) $\|\hat\alpha\psi\|^2 = \braket{\psi}{\hat\alpha^2\psi} =
\Delta A^2$ por la hermiticidad de $\hat\alpha$. (b) $\braket{\hat\alpha
\psi}{\hat\beta\psi} = \tfrac12\langle\{\hat\alpha,
\hat\beta\}\rangle + \tfrac12\langle[\hat A, \hat B]\rangle$: el
primer término es real (parte hermítica) y el segundo, imaginario puro.
(c) $|z|^2 \ge (\operatorname{Im}z)^2$ da el teorema; la igualdad
exige vectores proporcionales \emph{y} media nula del anticonmutador.
(d) Un estado propio de $\hat A$ tiene $\Delta A = 0$, lo que obliga a
$0 \ge \hbar/2$: imposible para estados normalizables. Las ondas planas
solo lo ``consiguen'' por ser idealizaciones no normalizables, fuera del
\omterm{def:b3:quantum-formalism:state}{espacio de Hilbert}.
\end{solution}

\begin{solution}{exo:b3:quantum-formalism:12}
(a) En $\omega t = \pi/2$, es decir, $t = T$. (b) En cada comprobación, el
estado ha girado $\pi/2N$; se lo halla en $\ket1$ con
$\cos^2(\pi/2N)$ y \emph{colapsa de vuelta} a $\ket1$; las comprobaciones
son independientes, de ahí el producto. (c) $0$, $0.53$, $0.88$ y $0.98$:
vigilada lo bastante de cerca, la olla nunca hierve. (d) El postulado de
proyección (P5): cada observación devuelve la evolución a su línea de
salida.
\end{solution}

\begin{solution}{pb:b3:quantum-formalism:1}
\textbf{1.} Una rotación lleva una pareja ortonormal a otra ortonormal:
$\braket{\nu_e}{\nu_\mu} = -\cos\theta\sin\theta +
\sin\theta\cos\theta = 0$.
\textbf{2.} P6: el vuelo libre lo genera $\hat H$, cuyos estados propios
evolucionan de forma autónoma mediante fases --- sea cual sea la base que
eligió la producción.
\textbf{3.} $\ket{\psi(0)} = \ket{\nu_\mu} = -\sin\theta\ket{\nu_1}
+ \cos\theta\ket{\nu_2}$.
\textbf{4.} $\ket{\psi(t)} = -\sin\theta\,\eu^{-\iu E_1t/\hbar}
\ket{\nu_1} + \cos\theta\,\eu^{-\iu E_2t/\hbar}\ket{\nu_2}$.
\textbf{5.} Las fases globales desaparecen de todo $|\braket\cdot\cdot|^2$:
queda $\Delta\phi = (E_2 - E_1)t/\hbar$.
\textbf{6.} $E_2 - E_1 = (m_2^2 - m_1^2)c^4/2E = \Delta m^2c^4/2E$.
\textbf{7.} $\braket{\nu_e}{\psi(t)} = \eu^{-\iu E_1t/\hbar}
\sin\theta\cos\theta\,(\eu^{-\iu\Delta\phi} - 1)$.
\textbf{8.} $|\cdots|^2 = \sin^2\theta\cos^2\theta\,|{\eu^{-\iu
\Delta\phi} - 1}|^2 = \sin^22\theta\,\sin^2(\Delta\phi/2)$.
\textbf{9.} Sin mezcla, o sin desdoblamiento de masas: no hay oscilación.
Toda oscilación observada demuestra $\theta \neq 0$ \emph{y} $m_1 \neq
m_2$ --- los neutrinos pesan.
\textbf{10.} Sustitúyase $\Delta\phi = \Delta m^2c^4L/2\hbar cE$;
$L_{\text{osc}}$ hace que el argumento valga $\pi$.
\textbf{11.} Las dos probabilidades de sabor son componentes al cuadrado en
una base ortonormal de un estado normalizado: suman $1$ ---
la unitariedad de la evolución conservó la norma.
\textbf{12.} Solo es \omterm{def:b3:quantum-formalism:observable}{observable} la fase \emph{relativa}, y contiene
$E_2 - E_1 \propto m_2^2 - m_1^2$: las masas absolutas se cancelan.
\textbf{13.} $\dfrac{L}{4\hbar cE}$ en las unidades enunciadas:
$\num{e3}\,\unit{m}/(4 \times \num{1.973e-7}\,\unit{eV.m} \times
\num{e9}) = 1.27$ por $\unit{eV^2}$.
\textbf{14.} Bien desarrollada desde abajo: $1.27\,\Delta m^2 \times 12800
\gtrsim 1$, es decir, $\Delta m^2c^4 \gtrsim \num{6e-5}\,\unit{eV^2}$;
despreciable desde arriba: $1.27\,\Delta m^2 \times 15 \ll 1$, es decir, $\ll
\num{5e-2}$: en torno a $10^{-4}$--$10^{-2}\,\unit{eV^2}$.
\textbf{15.} $L_{\text{osc}} = \pi E/(1.27\,\Delta m^2c^4) \approx
\qty{990}{km}$ a \qty{1}{GeV}: los \qty{15}{km} quedan intactos y los
\qty{12800}{km} son trece longitudes completas --- exactamente el patrón
observado.
\textbf{16.} A lo largo de muchas longitudes y de un abanico de energías,
$\langle\sin^2\rangle = \tfrac12$: supresión $\tfrac12\sin^22
\theta$; la mitad medida obliga a $\sin^22\theta \approx 1$ y
$\theta \approx 45^\circ$ --- la naturaleza eligió mezcla máxima.
\textbf{17.} $m_2c^2 \approx \sqrt{\num{2.5e-3}} = \qty{0.05}{eV}$:
diez millones de veces más ligero que el electrón --- la materia más ligera
conocida, y nadie sabe todavía por qué.
\textbf{18.} Daya Bay: $1.27 \times \num{2.5e-3} \times 1.5/0.004
\approx 1.2$ --- del orden de la unidad, sintonizado; KamLAND: $1.27 \times
\num{7.5e-5} \times 180/0.004 \approx 4.3$ --- del orden de la unidad para
el desdoblamiento solar: cada distancia es un interferómetro ajustado a un
$\Delta m^2$.
\textbf{19.} Problema: solo llegaba un tercio de los $\nu_e$ previstos del
Sol. Resolución: los dos tercios que faltan llegan como otros sabores, a
los que han rotado los $\nu_e$ (SNO contó el total y declaró inocente al
Sol).
\textbf{20.} La oscilación exige $\Delta m^2 \neq 0$: al menos un
neutrino tiene masa --- la primera física de laboratorio más allá del modelo
estándar original.
\textbf{21.} Mantener la coherencia de fase a lo largo de $10^7$ metros
significa que prácticamente \emph{nada} interaccionó por el camino:
secciones eficaces tan pequeñas que hacen falta kilotoneladas de agua para
capturar un puñado --- de ahí las cincuenta mil toneladas de
Super-Kamiokande.
\textbf{22.} Los operadores de sabor son diagonales en la base de sabor,
que no es la base de energía: $[\hat H, \text{sabor}] \neq 0$ ---
el sabor sencillamente no es una magnitud conservada del vuelo libre,
mientras que la energía y el momento sí lo son.
\textbf{23.} Los dos desdoblamientos difieren en un factor treinta: para un
$L/E$ dado, una oscilación está activa y la otra o congelada o
completamente promediada --- cada experimento ve un sistema efectivo de dos
niveles.
\textbf{24.} Base de producción $\neq$ base de propagación (la
rotación $\theta$); P6 suministra las dos fases; la \omterm{thm:b3:quantum-formalism:postulates}{regla de Born} convierte
la fase relativa en una probabilidad.
\textbf{25.} $\theta \approx 45^\circ$ y $\Delta m^2c^4 =
\qty{2.5e-3}{eV^2}$ dan a un neutrino muónico de GeV una longitud de
oscilación cercana a \qty{1000}{km}: álgebra lineal de dos niveles,
verificada a través del planeta, y un Premio Nobel por la desaparición de
medio flujo.
\end{solution}
